\begin{pastexam}[label={px:ncku-114-ee-ds-16}]{成大電機所 114 學年度 單選題第 16 題}
（單選；每題 5 分，答錯倒扣 2 分）

How many recursive calls would the call ``Fibonacci(10)'' generate in the following program?
\begin{minted}[fontsize=\small]{c}
int Fibonacci(int n) {
    if (n < 1) return 0;
    if (n < 3 ) return 1;
\end{minted}
a) 55 \quad b) 108 \quad c) 110 \quad d) 20 \quad e) None of the above
\end{pastexam}

\begin{pastsoln}
答案是 b) 108。

原卷的程式只印到第三行，缺了遞迴那一行和右大括號。 照題意補上
\texttt{return Fibonacci(n - 1) + Fibonacci(n - 2);}，這是費氏數列的標準遞迴寫法，$F(1) = F(2) = 1$。

設 $C(n)$ 是呼叫 \texttt{Fibonacci(n)} 時\textbf{總共}發生的呼叫次數（包含它自己）。 $n < 3$ 時直接
回傳，$C(n) = 1$；$n \ge 3$ 時自己一次，加上兩個子呼叫各自引發的次數：
\[
C(1) = C(2) = 1, \qquad C(n) = 1 + C(n - 1) + C(n - 2)\quad (n \ge 3).
\]
（資結書第~18.1.1~節就是用這種遞迴式數呼叫次數。）逐項算：
\begin{center}
\begin{tabular}{c*{10}{c}}
\toprule
$n$ & 1 & 2 & 3 & 4 & 5 & 6 & 7 & 8 & 9 & 10 \\
\midrule
$C(n)$ & 1 & 1 & 3 & 5 & 9 & 15 & 25 & 41 & 67 & 109 \\
$F(n)$ & 1 & 1 & 2 & 3 & 5 & 8 & 13 & 21 & 34 & 55 \\
\bottomrule
\end{tabular}
\end{center}
可以看出 $C(n) = 2F(n) - 1$：令 $D(n) = C(n) + 1$，遞迴式變成 $D(n) = D(n-1) + D(n-2)$，
$D(1) = D(2) = 2$，所以 $D(n) = 2F(n)$。

\texttt{Fibonacci(10)} 總共 $109$ 次呼叫，其中最外面那一次是題目的呼叫本身，它\textbf{產生}的遞迴
呼叫是 $109 - 1 = 108$ 次，選 b)。 a) $55$ 是 $F(10)$ 的值，不是呼叫次數；c) $110$ 多算了；d) $20$
太少。
\end{pastsoln}
